% =============================================================================
%                              MODULE OVERVIEW
% =============================================================================
\section*{Module Overview}
\addcontentsline{toc}{section}{Module Overview}

\subsection*{About This Module}

This module introduces students to \textbf{Natural Language Processing (NLP)} and the computational methods for processing human language, as well as the experimental methodology necessary for applying NLP in statistically valid settings.

\textbf{Module code:} 6CCSANLP \hfill \textbf{Semester:} 1 (AY 2025/26)

\bigseparator

\subsection*{Module Aims}

The module aims to:
\begin{itemize}
  \item Provide an introduction to NLP and cover computational methods for manipulating human language.
  \item Cover core tasks in NLP such as \textbf{sentiment analysis} and \textbf{machine translation}.
  \item Teach students how to \textbf{determine appropriate methods} for a range of NLP tasks.
  \item Develop skills to \textbf{implement} and \textbf{evaluate} NLP models.
\end{itemize}

\bigseparator

\subsection*{Learning Outcomes}

By the end of this module, you will be able to:

\begin{center}
\renewcommand{\arraystretch}{1.4}
\begin{tabular}{@{} c p{12cm} @{}}
\toprule
\textbf{Code} & \textbf{Learning Outcome} \\
\midrule
NLP1 & Understand the broad range of \textbf{methodologies and algorithms} related to language processing. \\
NLP2 & \textbf{Analyse and formulate} the objectives of different NLP problems. \\
NLP3 & Apply algorithms in \textbf{statistically valid experiments}. \\
NLP4 & Determine appropriate methods and know how to \textbf{implement them}. \\
NLP5 & Implement \textbf{machine learning algorithms} for NLP using a high-level programming language. \\
NLP6 & \textbf{Assess the success (or failure)} of an algorithm to learn a particular task. \\
\bottomrule
\end{tabular}
\end{center}

\bigseparator

